\documentclass[a4paper,
oneside,
12pt,
titlepage]
{scrartcl}   

\usepackage[sc]{mathpazo} % Schrift - wie Funcky und in PDF zu Fonts beschrieben
\usepackage[T1]{fontenc}
\usepackage[utf8]{inputenc}
\usepackage[a-1b]{pdfx}
\usepackage[ngerman]{babel}
\usepackage{textcomp} 
\usepackage[amssymb]{SIunits} 

\usepackage{graphicx} 
\usepackage{subfigure}                         
\usepackage{float}
\usepackage[section]{placeins}
\usepackage{commath} %Thecommath-package provides some commandswhich help you to format formulas flexibly
\usepackage[iso,german]{isodate} %his package provides commands to switch between different date formats
\usepackage{hyperref}
\usepackage{acronym}
\usepackage{listings}

\usepackage[backend=biber, style=alphabetic]{biblatex}
\usepackage{csquotes}
\addbibresource{Quellen.bib}


%Header Definitionen
\usepackage{fancyhdr}
\renewcommand{\headrulewidth}{0.5pt}
\renewcommand{\footrulewidth}{0.5pt}
%Abstand zwischen Absätzen, Zeilenabstände
\voffset26pt 
\parskip6pt
%\parindent1cm  %Rückt erste Zeile eines neuen Absatzes ein
\usepackage{setspace}
\onehalfspacing

%Hier beginnt das Dokument
\begin{document}
\pagenumbering{roman}

\titlehead
{
\begin{tabular}{ll}
\begin{minipage}{0.5\textwidth}
 \begin{figure}[H]
 \raggedright
  \includegraphics[scale=0.15]{bilder/MDT-logo}\\
\end{figure}
\end{minipage}
\begin{minipage}{0.5\textwidth}
\begin{figure}[H]
\raggedleft
\includegraphics[scale=0.04]{bilder/tu-logo}\\
\end{figure}
\end{minipage}
\end{tabular}\\
\\
  \small
      {
	Technische Universität Berlin\\
	Fakulät IV\\
	Institut für Energie und Automatisierungstechnik\\
	Fachgebiet Elektronische Mess- und Diagnosetechnik\\

	Praktikum Grundlagen der elektronischen Messtechnik\\
	Betreuer: Max Mustermann\\
	WS 2012/13}

}

\title {Versuchsprotokoll}
\subtitle{Versuch 1\\
Einführung in Python}
\author{Ein Kommilitone (123456)\\Niemand Besonderes (654321)}

\date{\today\\*[60pt]}   %Stellt das Datum ein /today wählt automatisch das aktuelle Datum, kann aber auch manuell eingestellt werden

\maketitle  %Erstellt das Titelblatt wie oben definiert

\thispagestyle{empty}

\newpage

%Einstellungen zur Kopf- und Fußzeile
\pagestyle{fancy}
\fancyhead[L]{Ein Kommilitone, Niemand Besonderes}	%Head = oben, foot = unten mit jeweiliger Option [L,C,R]
\fancyhead[R]{Praktikumsbericht: DKR}

\tableofcontents % Inhaltsverzeichnis, auskommentieren wenn unnötig
\newpage

\pagenumbering{arabic}

\section{Grundlagen der Parameteridentifikation linear, dynamischer Systeme}

...

\section{Regler}

...

\subsection{Allgemeine Struktur eines Regelkreises}

...

\subsection{PID-Regler}

...

\section{Regelung einer Drosselklappe}

...

\begin{figure}[H] %H Bild Hier, h Bild auf Seite
\centering
\includegraphics[width=0.8\textwidth]{bilder/tu-logo} 
\caption{Hier soll eine Abbildung hin.}
\label{fig:Versuchsaufbau} %Label idealerweise gleich wie Dateinamen (Zuordnung)
\end{figure}

...

\section{Drehmomentenbasierte Funktionsstruktur}

...

\begin{figure}[H] %H Bild Hier, h Bild auf Seite
\centering
\includegraphics[width=0.8\textwidth]{bilder/tu-logo} 
\caption{Hier soll eine Abbildung hin.}
\label{fig:Versuchsaufbau} %Label idealerweise gleich wie Dateinamen (Zuordnung)
\end{figure}

...

\subsection{Schneller und langsamer-Pfad}

...

\subsection{Sensoren und Aktoren}

...

\section{Ergebnisse der umgesetzten Funktionsstruktur}

...

\begin{figure}[H] %H Bild Hier, h Bild auf Seite
\centering
\includegraphics[width=0.8\textwidth]{bilder/tu-logo} 
\caption{Hier soll eine Abbildung hin.}
\label{fig:Versuchsaufbau} %Label idealerweise gleich wie Dateinamen (Zuordnung)
\end{figure}

...

\section{Aufbau einer digitalen Messketten zum erfassen analoger Signale}

...

\section{Ansteuerung von Akuatoren}

...

\subsection{High- und Low-Side-Schaltung}

...

\subsection{H-Brücke}

...

\subsection{PWM}

...

\printbibliography
\nocite{Kopka:Latex}

\end{document}
